\documentclass[12pt]{article}
\usepackage[T1]{fontenc}
\usepackage[utf8]{inputenc}
\usepackage{lmodern}
\usepackage{textcomp}
\usepackage{enumitem}
\usepackage{geometry}
\geometry{margin=1in}
\begin{document}
\begin{center}
Answer Key - Economics - K-12 Level Assessment 17 \
\small (Answers are ordered exactly as in the question paper.)
\end{center}

\section*{Multiple Choice}
\begin{enumerate}[label=\arabic*.]
\item \textbf{Answer:} C
\\ \textit{Options:} A) is ignored.; B) appears in the capital account.; C) appears in the current account; D) is included in the official reserves.
\\ \textit{Note:} The trade balance, which measures the difference between a country's exports and imports of goods and services, is a key component of the current account in the balance of payments.
\item \textbf{Answer:} C
\\ \textit{Options:} A) Increasing Increasing; B) Increasing Decreasing; C) Decreasing Increasing; D) Decreasing Unchanged
\\ \textit{Note:} Lower interest rates in the United States reduce the value of the dollar because they make U.S. investments less attractive to foreign investors, which in turn increases U.S. exports as American goods become cheaper for buyers in other countries.
\item \textbf{Answer:} B
\\ \textit{Options:} A) an increase in the real interest rate; B) increased investment in capital; C) an increase in aggregate demand; D) an increase in the unemployment rate
\\ \textit{Note:} Increased investment in capital enhances productivity and the capacity of the economy to produce goods and services, leading to an increase in long-run aggregate supply.
\item \textbf{Answer:} A
\\ \textit{Options:} A) diminishing marginal utility; B) diminishing marginal returns; C) the Fisher effect; D) diminishing returns to scale
\\ \textit{Note:} The demand curve slopes downward for an individual because as they consume more units of a good, the additional satisfaction or utility gained from each additional unit decreases, leading them to be willing to pay less for each subsequent unit.
\item \textbf{Answer:} A
\\ \textit{Options:} A) A decrease in taxes and a lower discount rate; B) An increase in government spending and an increase in taxes; C) A decrease in taxes and selling bonds in an open market operation; D) An increase in government spending and an increase in the discount rate
\\ \textit{Note:} A decrease in taxes increases disposable income for consumers and businesses, stimulating demand and economic activity, while a lower discount rate reduces borrowing costs, encouraging investment, both contributing to a rise in real GDP above full employment levels.
\end{enumerate}

\section*{True / False}
\begin{enumerate}[label=\arabic*.]
\setcounter{enumi}{5}
\item \textbf{Answer:} True
\\ \textit{Options:} A) True; B) False
\\ \textit{Note:} An increase in consumer confidence typically leads to higher consumer spending, which raises demand for goods and services, resulting in an increase in both the equilibrium price level and the equilibrium quantity of output in the economy.
\item \textbf{Answer:} False
\\ \textit{Options:} A) True; B) False
\\ \textit{Note:} The statement is false because a marginal propensity to save (MPS) of 0.20 means that for every $100 increase in income, saving increases by $20, but the statement incorrectly suggests that this is a correct interpretation without acknowledging that it might be misleading in context.
\item \textbf{Answer:} True
\\ \textit{Options:} A) True; B) False
\\ \textit{Note:} The Phillips curve demonstrates that as inflation rates rise, unemployment rates tend to decrease, indicating an inverse relationship between the two economic indicators.
\item \textbf{Answer:} True
\\ \textit{Options:} A) True; B) False
\\ \textit{Note:} The statement is true because allocative efficiency occurs when resources are distributed in a way that maximizes consumer satisfaction, which is achieved when the price consumers are willing to pay for a good or service equals the additional cost of producing one more unit of that good or service (marginal cost).
\item \textbf{Answer:} True
\\ \textit{Options:} A) True; B) False
\\ \textit{Note:} A balance of payments deficit indicates that a country is spending more on foreign trade than it is earning, leading to a depletion of its official reserves to cover the difference.
\end{enumerate}

\section*{Long Form Response}
\begin{enumerate}[label=\arabic*.]
\setcounter{enumi}{10}
\item \textbf{Answer:} The ownership of a unique item, like the only original signed copy of The Wealth of Nations by Adam Smith, can be represented by a perfectly vertical supply curve in economics because there is only one of that item available in the market. A perfectly vertical supply curve indicates that no matter the price, the quantity supplied remains constant at one unit. This reflects the fact that since there is only one original signed copy, its supply cannot increase regardless of demand. Therefore, the price may fluctuate based on how much buyers are willing to pay, but the quantity remains fixed, illustrating the concept of perfectly inelastic supply.
\\ \textit{Note:} correct because a perfectly vertical supply curve demonstrates that the supply of a unique item, such as the only original signed copy of The Wealth of Nations, is fixed at one unit, meaning its quantity cannot change in response to price fluctuations.
\item \textbf{Answer:} Inflation at the peak of a business cycle can threaten the macroeconomy by reducing consumers' purchasing power, increasing costs for businesses, and leading to overall economic instability. When prices rise rapidly, consumers find that their money buys less, which can lead to decreased spending and lower demand for goods and services. For businesses, higher costs for materials and labor can squeeze profits, prompting them to raise prices even further or cut back on investment and hiring. This cycle can create uncertainty in the economy, discouraging both consumer and business confidence, which may lead to a slowdown or recession if not addressed. Ultimately, sustained inflation can disrupt economic growth and create challenges for policymakers trying to maintain stability.
\\ \textit{Note:} correct because it identifies the negative impacts of inflation during a business cycle peak, including diminished consumer purchasing power, increased business costs, and potential economic instability, all of which can contribute to a downward economic spiral.
\end{enumerate}

\vfill
\noindent\textit{End of Answer Key}
\end{document}